\documentclass[10pt,a4paper]{amsart}
\usepackage[margin=19mm]{geometry}
\usepackage{amsmath,amssymb,mathtools}
\usepackage[scr=rsfs,cal=euler]{mathalfa}
\usepackage{xfrac}
\usepackage[unicode,hidelinks,bookmarks=false]{hyperref}
\newcommand{\ie}{i.e.\ }
\newcommand{\cf}{cf.\ }
\newcommand{\resp}{respectively\ }
\newcommand{\Subsection}{\subsection}
\providecommand{\nobreakdash}{\nobreak\hbox{-}}
\setlength{\emergencystretch}{2em}
\multlinegap0pt
\usepackage{amssymb}
\usepackage{xcolor}
\usepackage{float}
\usepackage{algorithm}\usepackage{algpseudocode}
\usepackage{tikz}
\usepackage{mathtools}
\newtheorem{lemma}{Lemma}
\title{INTHOP: A Second-Order Globally Convergent Method for Nonconvex Optimization}
\author{Krishan Kumar, Ashutosh Sharma, Gauransh Dingwani, Nikhil Gupta, Vaishnavi Gupta, Ishan Bajaj}
\date{Source: arXiv:2510.22342v4 (CC BY 4.0)}
\begin{document}
\maketitle

\noindent\emph{Source and licence.} INTHOP: A Second-Order Globally Convergent Method for Nonconvex Optimization. Version arXiv:2510.22342v4 (CC BY 4.0), distributed by arXiv under the Creative Commons Attribution 4.0 International licence, http://creativecommons.org/licenses/by/4.0/, as recorded in the arXiv licence metadata for this exact version and shown on the abstract page https://arxiv.org/abs/2510.22342v4. Full licence notice: \texttt{licenses/arxiv-2510.22342-CC-BY-4.0.txt}.\par\smallskip

\noindent\textit{Editorial scope.} Counted unit: the article's lemma steplengthlemma (source0.tex lines 771--776). The article's complete original proof of it is delivered with it (source0.tex lines 777--816, one \texttt{proof} environment of 40 lines). No mathematics is written, repaired or re-derived: the statement and the proof are the article's own lines, in the article's own notation, and no retained line is substituted or re-flowed.  Retained uncounted as the source-local context the proof needs: lines 438--440; lines 764--766.  Nothing was omitted from any retained span, so the package carries no omission record.  No retained line contains a citation, so the wrapper carries no bibliography and no bibliography entry.  No retained line contains a figure environment, an image-inclusion command or drawing code, so no image is delivered and the package declares no asset dependency.  Editorial formatting only, in addition: an AMS \texttt{amsart} preamble that restates the article's own declarations and macros used by the retained lines, namely the environments \texttt{lemma} on separate counters, together with the standard packages those lines need.  The retained environments print in this document as Lemma 1 in order of appearance, whereas the article prints the same items with the numbering of its own full document.  The complete native source of the article as distributed in the arXiv e-print for this exact version is bundled unchanged as \texttt{sources/187786/source0.tex}.
    \begin{equation}\label{p_k}
        \mathbf{p}_k = - \nabla^2 \mathcal{H}_t^{-1}\mathbf{g}_k
    \end{equation}

    \begin{equation}\label{steplengthpf3}
        \lVert \nabla^{2} f(x)\rVert \leq \mathcal{A} \text{ for all } x\in \mathcal{S}+\mathbb{B}(0,\mathcal{D}).
    \end{equation}

\begin{lemma}\label{steplengthlemma}
    Suppose that $\theta_{k}$ is the step length generated by the algorithm using Armijo line search condition for each $k\geq0$. Then,  
    \begin{equation*}
    \theta_{k}\geq \min\left(1,\frac{2(1-\eta)c_1\tilde{g}}{\mathcal{A}}{t}\right),  \text{ where } 0<t<1 \text{ and } \mathcal{A} \text{ is given in }\eqref{steplengthpf3}.      
    \end{equation*}
\end{lemma}

\begin{proof}
    Note that $f$ is twice differentiable function. Therefore, from Taylor's theorem, there exists a scalar $b\in (0,1)$ such that 
    \begin{align}\label{steplengthpf1}
    &f(\mathbf{x}_{k}+\theta_{k}\mathbf{p}_k)=-f(\mathbf{x}_{k})+\theta_{k}\mathbf{g}^{\top}_{k}\mathbf{p}_{k}+\frac{1}{2} \theta^{2}_{k}\mathbf{p}^{\top}_{k} \nabla^{2} f(\mathbf{x}_{k}+b\theta_{k}\mathbf{p}_{k})\mathbf{p}_{k}\nonumber\\
    \implies& f(\mathbf{x}_{k})-f(\mathbf{x}_{k}+\theta_{k}\mathbf{p}_{k})+\eta\theta_{k}\mathbf{g}^{\top}_{k}\mathbf{p}_{k}=-(1-\eta)\theta_{k}\mathbf{g}^{\top}_{k}\mathbf{p}_{k}-\frac{1}{2} \theta^{2}_{k}\mathbf{p}^{\top}_{k}\nabla^{2}f(\mathbf{x}_{k}+b\theta_{k}\mathbf{p}_{k})\mathbf{p}_{k}.
    \end{align}
    From \eqref{p_k}, we have $\mathbf{p}_{k}=-\nabla^{2}\mathcal{H}^{-1}_{t}\mathbf{g}_k$, where  $\nabla^{2}\mathcal{H}^{-1}_{t}=(\nabla^{2} f(\mathbf{x}_{t})+2\alpha I+c_1\tilde g I)^{-1}$ and $\mathbf{x}_{t}\in
    [\mathbf{x}^{L}_{t}, \mathbf{x}^{U}_{t}]$. Therefore,
    \begin{equation}\label{steplengthpf2}
    \mathbf{g}_k=-(\nabla^{2} f(\mathbf{x}_{t})+2\alpha I+c_1\tilde g I)\mathbf{p}_{k}.
    \end{equation}
    In view of \eqref{steplengthpf1} and \eqref{steplengthpf2}, we get 
    \begin{align}\label{steplengthpf4}
        &~~~~f(\mathbf{x}_{k})-f(\mathbf{x}_{k}+\theta_{k}\mathbf{p}_{k})+\eta\theta_{k}\mathbf{g}^{\top}_{k}d_{k}\nonumber\\
        &= (1-\eta) \theta_{k} \mathbf{p}^{\top}_{k} (\nabla^{2} f(\mathbf{x}_{t})+2\alpha I+c_1\tilde g I)\mathbf{p}_{k}-\frac{1}{2} \theta^{2}_{k} \mathbf{p}^{\top}_{k}\nabla^{2} f(\mathbf{x}_{k}+b\theta_{k}\mathbf{p}_{k})\mathbf{p}_k\nonumber\\
        &= (1-\eta)\theta_{k}\mathbf{p}^{\top}_{k}(\nabla^{2} f(\mathbf{x}_{t})+2\alpha I)\mathbf{p}_{k}+(1-\eta)\theta_{k}\mathbf{p}^{\top}_{k}\left(c_1\tilde gI-\frac{1}{2(1-\eta)}\theta_{k}\nabla^{2} f(\mathbf{x}_{k}+b\theta_{k}\mathbf{p}_{k})\right)\mathbf{p}_{k}\nonumber\\
        &\geq (1-\eta) \theta_{k} \mathbf{p}^{\top}_{k} \left(c_1\tilde gI-\frac{1}{2(1-\eta)}\theta_{k}\nabla^{2} f(\mathbf{x}_{k}+b\theta_{k}\mathbf{p}_{k})\right)\mathbf{p}_{k}\nonumber\\
        &~~~~~~~~~~~~~~~~~~~~~~~~~~~~~~~~~~~~~~~~~~~~~~~\text{ as }\nabla^{2} f(\mathbf{x}_{t})+2\alpha I \text{ is positive semidefinite  for }\mathbf{x}_{t}\in[\mathbf{x}^{L}_{t}, \mathbf{x}^{U}_{t}]\nonumber\\
        &=(1-\eta)\theta_{k}\left(c_1\tilde g\lVert \mathbf{p}_{k}\rVert^{2}I-\frac{1}{2(1-\eta)}\theta_{k}\mathbf{p}^{\top}_{k}\nabla^{2}f(\mathbf{x}_{k}+b\theta_{k}\mathbf{p}_{k})\mathbf{p}_{k}\right)\nonumber\\
        &\geq (1-\eta)\theta_{k}\left(c_1\tilde g-\frac{1}{2(1-\eta)}\theta_{k}\rVert\nabla^{2}f(\mathbf{x}_{k}+b\theta_{k}\mathbf{p}_{k})\rVert\right)\lVert \mathbf{p}_{k}\rVert^{2}.
    \end{align}
    
    Now, from  \eqref{steplengthpf4}, and \eqref{steplengthpf3}, we get 
    \begin{align}\label{steplengthpf5}
        f(\mathbf{x}_{k})-f(\mathbf{x}_{k}+\theta_{k}\mathbf{p}_{k})+\eta\theta_{k}\mathbf{g}^{\top}_{k}d_{k}\geq (1-\eta) \theta_{k} \left(c_1\tilde{g}-\frac{1}{2(1-\eta)}\theta_{k}\mathcal{A}\right)\lVert \mathbf{p}_{k}\rVert^{2}.
    \end{align}
Next, we consider two cases.\\
First, if $\frac{2(1-\eta)c_1\tilde{g}}{\mathcal{A}}\geq1$, then from \eqref{steplengthpf5}, 
\begin{equation*}
f(\mathbf{x}_{k})-f(\mathbf{x}_{k}+\mathbf{p}_k)\geq-\eta \mathbf{g}^{\top}_{k}\mathbf{p}_{k}
\end{equation*}
which implies that $\theta_{k}=1$ satisfies Armijo line search condition.\\
Second, if $\frac{2(1-\eta)c_1\tilde{g}}{\mathcal{A}}<1$, then for $\theta_{k}\leq \frac{2(1-\eta)c_1\tilde{g}}{\mathcal{A}}$ in \eqref{steplengthpf5}, we have 
\begin{equation}
    f(\mathbf{x}_{k})-f(\mathbf{x}_{k}+\mathbf{p}_k)\geq-\eta \mathbf{g}^{\top}_{k}\mathbf{p}_{k}
\end{equation}
which implies that $\theta_{k}$ must be greater than or equal to $\frac{2(1-\eta)c_1\tilde{g}}{\mathcal{A}}{t_1}$ for $0<t_1<1$. Otherwise, there will be a $\tilde{\theta_{k}}=\frac{\theta_{k}}{t_1}>\theta_{k}$ that 
satisfies 
$ f(\mathbf{x}_{k})-f(\mathbf{x}_{k}+\tilde{\theta_{k}}\mathbf{p}_k)\geq-\eta\tilde{\theta_{k}} \mathbf{g}^{\top}_{k}\mathbf{p}_{k}$ which contradicts the definition of $\theta_{k}$. Hence, we have $\theta_{k}\geq\min\left(1,\frac{2(1-\eta)c_1\tilde{g}}{\mathcal{A}}{t_1}\right)$, where  $0<t_1<1$. 
\end{proof}

\end{document}
